\chapter{Development Process and Method}\label{ch:method}

The project follows Scrum \cite{schwaber-scrum-guide-2016} in two-week sprints. This
chapter describes how the group has set up roles, sprints and the product backlog,
where the setup deviates from the Scrum Guide, and how progress is measured. The
practical rules are in Appendix 10.1, and the decisions are logged with dates in
Appendix 10.2.

\section{Roles}\label{sec:method-roles}

Frederik is Scrum Master. The rest of the group is the development team.

The Scrum Guide describes the Product Owner as the one person who orders the product
backlog: ``The Product Owner is one person, not a committee''
\cite[p.~5]{schwaber-scrum-guide-2016}. We deviate from that. Our Product Owner is
Bodil, a blue toy elephant, and the whole group orders the backlog together at
sprint planning. Bodil is the one the group talks to when a feature is discussed: the
question is what Bodil, as a customer, would need. The idea came as advice from a
previous exam.

The cost of the deviation is that no single person settles a disagreement about
priority. The MoSCoW priorities in the requirements specification (Appendix 1.2) are
therefore the tie-breaker: a MUST requirement goes before a SHOULD, and a task that
does not trace to a requirement has to be argued for.

\section{Sprints}\label{sec:method-sprints}

A sprint runs from a Wednesday to the Tuesday two weeks later. Sprint 1 ran from 16 to
29 September, and sprint 6 ends on 8 December, three days before the report is handed
in. A sprint starts with sprint planning, where the group picks the tasks for the
sprint and writes a sprint goal, and ends with a sprint review and a retrospective.
The sprint goal is one sentence that names the requirements the sprint delivers, so
the tasks in a sprint belong together \cite[p.~10]{schwaber-scrum-guide-2016}. The
last sprint is set aside for the full acceptance test and the report.

\section{Product Backlog and Board}\label{sec:method-backlog}

The backlog lives in a self-hosted Kaneo instance, which replaced Notion on
9 September. Kaneo is an open-source alternative to Trello. It has what we need, a
backlog, a board, labels and due dates, without the extra fields and configuration of
larger tools such as Jira, so the board stays easy to read for everyone in the group.
It can also be read by the AI tool the group uses.
Kaneo's Backlog view is the product backlog, and the board with the columns
\emph{To Do}, \emph{In Progress}, \emph{In Review} and \emph{Done} is the sprint
backlog.

A task moves from the backlog to \emph{To Do} only at sprint planning, and only when
it is ready: its description starts with the requirements it covers, it has a
definition of done, a story point estimate, an epic label and an owner. The Scrum
Guide calls an item ready when it can be done within one sprint
\cite[p.~14]{schwaber-scrum-guide-2016}; the list above is how we check that. The epic
labels follow the grouping of the user stories in the requirements specification, so
every task can be traced back to a requirement.

The AI tool may create tasks and suggest estimates, but moving, assigning and
estimating tasks is done by the group. Planning a sprint is part of what the project
is meant to teach us, so it is not handed to a tool.

\section{Estimation and Progress}\label{sec:method-progress}

Tasks are estimated in story points with planning poker, on the scale 0, \textonehalf,
1, 2, 3, 5, 8, 13, 20, 40 and ``?'' for a task that cannot be estimated yet. As the
Scrum Guide prescribes, the people who do the work make the estimate
\cite[p.~14]{schwaber-scrum-guide-2016}.

At sprint planning we note the points committed to the sprint, and at the sprint
review the points done. Together with the size of the whole backlog, the numbers are
kept in one table (Appendix 10.2), which Figure~\ref{fig:burnup} is drawn from. We use
a burn-up rather than a burn-down chart. Most of the backlog was not estimated when
the project started, so the total will grow, and a burn-down chart would show that
growth as missing progress. A burn-up chart keeps the two apart
\cite[p.~14]{schwaber-scrum-guide-2016}.

\begin{figure}[H]
    \centering
    \pgfplotstableread[col sep=comma]{assets/data/sprint-velocity.csv}\velocitytable
    \def\velocitydone{0}
    \pgfplotstablecreatecol[create col/assign/.code={%
        \getthisrow{done}\entry
        \ifx\entry\empty
            \pgfkeyssetvalue{/pgfplots/table/create col/next content}{nan}%
        \else
            \pgfmathparse{\velocitydone+\entry}\xdef\velocitydone{\pgfmathresult}%
            \pgfkeyslet{/pgfplots/table/create col/next content}\velocitydone
        \fi}]{donetotal}\velocitytable
    \pgfplotstablecreatecol[create col/assign/.code={%
        \getthisrow{scope}\entry
        \ifx\entry\empty
            \pgfkeyssetvalue{/pgfplots/table/create col/next content}{nan}%
        \else
            \pgfkeyslet{/pgfplots/table/create col/next content}\entry
        \fi}]{scopetotal}\velocitytable
    \begin{tikzpicture}
        \begin{axis}[
            width=0.8\textwidth, height=6cm,
            xlabel={Sprint}, ylabel={Story points},
            xmin=1, xmax=6, xtick={1,...,6}, ymin=0,
            legend pos=south east, unbounded coords=jump,
        ]
            \addlegendimage{blue, mark=*}
            \addlegendentry{Backlog size}
            \addlegendimage{red, mark=square*}
            \addlegendentry{Done (cumulative)}
            \addplot[draw=none, forget plot] table[x=sprint, y expr={2*(\thisrow{sprint}-1)}]{\velocitytable};
            \addplot[blue, mark=*, forget plot] table[x=sprint, y=scopetotal]{\velocitytable};
            \addplot[red, mark=square*, forget plot] table[x=sprint, y=donetotal]{\velocitytable};
        \end{axis}
    \end{tikzpicture}
    \caption[Burn-up per sprint]{Burn-up per sprint. The gap between the lines is the work left; when the
    upper line rises, the backlog has grown, not the progress slowed.}
    \label{fig:burnup}
\end{figure}
